\section{Related Work}

\subsection{Efficient Visual Token Reduction}

The quadratic complexity of self-attention with respect to sequence length has motivated a
substantial body of work on reducing the number of visual tokens processed by vision transformers.
Early approaches learn to \emph{prune} uninformative tokens during inference.
DynamicViT~\cite{rao2021dynamicvit} trains a lightweight prediction network to hierarchically
discard tokens at multiple stages of a ViT, guided by a differentiable token sparsification loss.
EViT~\cite{liang2022evit} identifies inattentive tokens via attention scores to the
\texttt{[CLS]} token and fuses them into a single representative token rather than discarding
them outright, preserving global context.
A-ViT~\cite{yin2022avit} extends this idea by introducing adaptive token halting with an
ACT-style~\cite{graves2016act} mechanism, halting each token's computation individually based
on its predicted informativeness.

A complementary strategy is \emph{token merging}, which combines redundant tokens rather than
removing them.
ToMe~\cite{bolya2023tome} greedily matches and averages the most similar token pairs via bipartite
soft matching, achieving significant throughput gains with minimal accuracy loss on image
classification and video tasks.
Our bipartite token merging in Section~\ref{sec:merge} directly builds on this mechanism,
replacing the similarity-based matching criterion with trajectory-guided importance scores so that
merging decisions are informed by egocentric behavioral context.

\subsection{Efficient Video Understanding with Large Language Models}

Adapting large language models (LLMs) to video requires compressing temporally extended visual
sequences into a token budget compatible with the LLM's context window.
Video-LLaVA~\cite{lin2023videollava} projects visual frames into a unified representation space
before feeding them to an LLM, but uses dense per-frame tokens that scale poorly with clip length.
LLaMA-VID~\cite{li2023llamavid} addresses this by representing each frame with only two tokens—a
context token summarizing global content and a content token encoding local detail—enabling
hour-long video understanding within fixed context limits.
Chat-UniVi~\cite{jin2024chatunivi} clusters visual tokens dynamically across frames, merging
spatiotemporally similar regions to reduce redundancy.
Qwen2-VL~\cite{wang2024qwen2vl}, our base model, employs a Naive Dynamic Resolution mechanism and
temporal compression via 3D convolution to control token count, but still relies on uniform
spatial coverage rather than semantic salience.

These works reduce visual tokens through generic compression strategies—uniform subsampling,
clustering, or resolution scaling—without access to behavioral priors about where a human agent
directs attention.
TrajGazeMerge is complementary: rather than compressing visual representations post-hoc, we use
egocentric gaze and hand trajectories to \emph{selectively preserve} the tokens most likely to
contain task-relevant information before the VLM even processes them.

\subsection{Egocentric Video Understanding}

Egocentric video understanding has been advanced by large-scale datasets capturing diverse
first-person activities.
Ego4D~\cite{grauman2022ego4d} provides 3,670 hours of egocentric footage annotated for
episodic memory, forecasting, and social interaction tasks.
EgoExo4D~\cite{grauman2024egoexo4d} extends this to paired egocentric and exocentric views of
skilled activities, enabling cross-view learning.
EPIC-Kitchens~\cite{damen2022rescaling} offers densely annotated kitchen activity videos with
fine-grained action labels.
EGTEA~\cite{li2018egtea} is notable for providing simultaneous eye-tracking data alongside
activity annotations, directly linking gaze behavior to action recognition.
HoloAssist~\cite{wang2023holoassist} focuses on procedural assistance, recording participants
performing tasks while receiving real-time guidance.

StreamGaze\_v2, the benchmark used in this work, is built upon EGTEA, EgoExo4D, and HoloAssist
clips and targets multiple-choice video QA requiring spatial, temporal, and gaze-predictive
reasoning.

Vision-language pretraining on egocentric data has emerged as a powerful paradigm for
learning grounded representations.
EgoVLP~\cite{lin2022egovlp} performs video-text contrastive pretraining on Ego4D, producing
egocentric representations that transfer well to downstream tasks including action recognition
and natural language query localization.
EgoVLPv2~\cite{pramanick2023egovlpv2} improves upon this by fusing cross-modal attention
directly into the video and text backbones rather than applying it as a separate late-fusion step.
EgoSchema~\cite{mangalam2024egoschema} introduces a challenging long-form VideoQA benchmark
that requires multi-minute temporal reasoning over egocentric clips, exposing limitations of
short-clip VLMs.

\subsection{Gaze and Hand Attention in Egocentric Video}

Gaze has long been recognized as a privileged supervisory signal in egocentric settings.
Li et al.~\cite{li2018egtea} demonstrate that joint modeling of gaze and hand action improves
first-person activity recognition, as gaze fixations strongly correlate with objects being
manipulated.
Huang et al.~\cite{huang2018predicting} predict future gaze positions from past context using a
task-dependent attention transition model, showing that gaze is structured and predictable from
activity context.
Nagarajan et al.~\cite{nagarajan2020ego} use egocentric gaze and motion to discover affordance
hotspots in the environment, demonstrating that gaze reveals task-relevant scene regions without
explicit annotation.

These works treat gaze as an \emph{output} to be predicted or a supervisory signal for
learning visual representations.
In contrast, TrajGazeMerge treats recorded gaze and hand trajectories as \emph{inputs} that
actively guide visual token selection inside a VLM, directly connecting egocentric behavioral
signals to efficient video language modeling.
